\begin{papercard}{duguay2023}
\cf{Stages} \stagetag{D} \stagetag{W} (also \stagetag{H}). A critical update of the Light et al.\ walkthrough method by one of its authors. It reviews later uses, reports teaching experience, and runs a TikTok walkthrough to show where the method fails: data flows, algorithms, personalisation, updates and users.

\cf{Aim and design} Explores the method's ``enduring strengths and emergent weaknesses'' against developments in app studies, with TikTok as the central example \pc{duguay2023}{1}{This article explores the method’s enduring strengths and emergent weaknesses}. Three sources of evidence: a literature review of adjacent methods, lessons from teaching (from workshops to undergraduate courses), and a technical walkthrough of TikTok on iOS in May 2022 by the second author \pc{duguay2023}{6}{Gold-Apel conducted a technical walkthrough in May 2022 with the present Apple iOS version of TikTok}. The original method deliberately leaves the theoretical approach and order of procedure open \pc{duguay2023}{1}{By leaving open certain considerations}.

\cf{Gaps and adjacent methods} \emph{Platform types}: the method treats all apps alike, but sectoral, connector and infrastructural platforms differ, e.g.\ Yang et al.'s ``backend-in'' approach \pc{duguay2023}{2}{While the walkthrough method made no distinction between these different types of platforms}. \emph{Data flows}: field notes can record registration data and permissions, and the policy or terms often reveal the business model, but the walkthrough ``fails to capture data processes that are not made explicit'' \pc{duguay2023}{3}{The walkthrough method can be attuned to data exchanges between users and apps}, \pc{duguay2023}{3}{However, the walk- through fails to capture data processes that are not made explicit}. Supplements: network sniffing and packet inspection, or Exodus with legal and ethical frameworks (Kuntsman et al.'s data-justice toolkit) \pc{duguay2023}{3}{Incorporating the walkthrough method into an approach that uses such detec}. \emph{Research personas}: a ``clean'' account may not represent everyday use, so personas should be designed deliberately \pc{duguay2023}{3}{weighing up whether a “clean” account with few data inputs}. \emph{History}: the walkthrough is a snapshot. Platform biography and platform historiography (archived SDK and API documentation) add a longitudinal view \pc{duguay2023}{3}{Calls for iterative explorations add a longitudinal compo}.

\cf{Can users be bracketed?} Users are caught up anyway, through interaction with research accounts and their data in field notes and screenshots \pc{duguay2023}{4}{users are inevitably caught up in the research pro}. Feminist data studies call for reflexivity and an ethics of care, and treating users only as props risks reproducing extractive logics \pc{duguay2023}{4}{When users are treated only as props or exam}. Adjacent participant-led methods are the scrollback method and the media go-along. Ethical conduct should cover collection, storage and de-identification of screenshots, and also what is lost by bracketing users \pc{duguay2023}{4}{how to deidentify users in screenshots and exam}.

\cf{Teaching lessons} In the ``comparative app walkthrough'', small groups compare two apps of the same genre, an inductive first step towards research questions \pc{duguay2023}{4}{small groups of learners chose two apps from the same genre}. Example: Facebook's guilt-laden deletion flow vs.\ VKontakte's abrupt one \pc{duguay2023}{5}{the difference between how Facebook’s account deletion functionality guilts users}. Connector apps produced null findings because the end user is not the key user \pc{duguay2023}{5}{The apps were infrastructural platforms and com}. Students were wary of sharing data, avoided single sign-on, worried about forgetting to delete apps, and personal data leaked into screenshots \pc{duguay2023}{5}{Students worried about forgetting to delete the app later}. The method assumes access to a suitable device and safe boundaries between personal and research use, which is not true for everyone \pc{duguay2023}{5}{the method assumes learners and researchers have access to the appropriate technical resources}. Some features are only visible through interaction, so some learners staged interactions among group members \pc{duguay2023}{5}{Some learners simply innovated on the method}.

\cf{TikTok stumbling blocks} (1) \emph{Centrality of the For You feed}: the algorithmic feed is central, but the method has no procedure for algorithms. Audits and reverse engineering need rigorous logs of inputs and outputs \pc{duguay2023}{6}{rigorously recording the actions they take in app and documenting the algorithm’s outputs}. (2) \emph{Hyper-individualisation}: registration via a new e-mail, birthday-based forks (users under 16 get private accounts), interest prompts and contact access \pc{duguay2023}{7}{she registered with an email address newly created for this purpose}. The researcher's own actions continuously alter the app, and automated collection is imperceptible \pc{duguay2023}{7}{such data are compiled in conjunction with volumes of automated data collection}. Zulli and Zulli ran two contrasting personas, but even an empty control account is not neutral \pc{duguay2023}{7}{even a relatively empty account, intended to act as a control or baseline}. (3) \emph{Ephemeral interfaces}: an update replaced Discover with Friends during the walkthrough and was reversed three months later. One walkthrough could not judge its impact \pc{duguay2023}{7}{While Gold-Apel recorded this change in field notes during the technical walk}. Interfaces vary by location, age and privacy jurisdiction \pc{duguay2023}{8}{depending on location, age, privacy jurisdiction}.

\cf{Conclusion} Keep the walkthrough as a ``springboard'' and supplement, extend or interchange it with other methods, rather than building one all-encompassing method \pc{duguay2023}{8}{There is no sense in trying to subsume all approaches to the study of apps under one comprehensive method}.

\cf{Use for this thesis} Methodological cautions for stage D/W. (1) Record the \emph{research persona and account configuration}: a new e-mail, no single sign-on, declared age and region, and an isolated device. Treat even a clean account as non-neutral. (2) Record the app version, OS and date, and note interface changes during data collection. Jurisdiction-dependent interfaces matter for a Germany-specific sample. (3) State explicitly that the walkthrough captures only \emph{user-visible} data practices, not hidden flows, so claims stay within what users can see. (4) De-identify any third-party content in screenshots, and do not use a personal device or account.
\end{papercard}
